\section*{Chapter \ref{ch:rs:market-data-for-research} --- Market Data for Research}

\begin{solution}{exo:rs:market-data-for-research:1}
$(300 \times 20.00 + 500 \times 20.02 + 200 \times 19.99)/1\,000 = 20\,008/1\,000 = 20.008$.
\end{solution}

\begin{solution}{exo:rs:market-data-for-research:2}
The spread is $0.01/20 = 0.0005$ in log terms, so each trade-to-trade return gains $0.0005^2/2 = 1.25 \times 10^{-7}$ of variance;
over 20\,000 returns that is $0.0025$. With the true $0.02^2 = 0.0004$ the sum is $0.0029$: a reported volatility of 5.39\%, 2.7
times the truth.
\end{solution}

\begin{solution}{exo:rs:market-data-for-research:3}
$40 \times 1.5 \times 10^6 \times 5\,000 \times 252 = 7.56 \times 10^{13}$ bytes, about 76 terabytes a year before compression.
\end{solution}

\begin{solution}{exo:rs:market-data-for-research:4}
$3(1 + 0.8^2) = 4.92$. A kurtosis of 6 needs $1 + \mathrm{CV}^2 = 2$, a coefficient of variation of 1.
\end{solution}

\begin{solution}{exo:rs:market-data-for-research:5}
Held: 70, 71, 76 (the next contract from the third day). The roll gap on day two is $74 - 71 = 3$. Back-adjusted: 73, 74, 76.
Ratio-adjusted: $70 \times 74/71 = 72.96$, $74$, $76$. The held position's return on day three is $76/74 - 1 = 2.7\%$, which the
ratio-adjusted series shows and the raw nearby series ($72/71$ then a jump) does not.
\end{solution}

\begin{solution}{exo:rs:market-data-for-research:6}
The same traded value is ten times as many shares at USD~10, so a threshold in shares closes ten times as many \omterm{def:rs:market-data-for-research:bar}{bars} a day in
2015 as in 2025. \omterm{def:rs:market-data-for-research:bar}{Dollar bars} close the same number in both years: value, not shares, measures the amount of trading.
\end{solution}

\begin{solution}{exo:rs:market-data-for-research:7}
3\,472 \omterm{def:rs:market-data-for-research:bar}{bars} of 16.0 trades on average (the median \omterm{def:rs:market-data-for-research:bar}{bar} has 7), with a midquote-return kurtosis of 35.4 (24.2 close to close).
The threshold is re-estimated from the \omterm{def:rs:market-data-for-research:bar}{bars} just closed. Because order signs are autocorrelated (metaorders), the imbalance
reaches $\sqrt{\E[T]}$ in fewer than $\E[T]$ trades, the estimate of $\E[T]$ falls, and the threshold with it: the \omterm{def:rs:market-data-for-research:bar}{bars}
shrink until they close on short runs of one-sided flow, which is where large mid moves happen.
\end{solution}

\begin{solution}{exo:rs:market-data-for-research:8}
Both signals use levels of a back-adjusted series, which are not prices: each has had later roll gaps added, so the ratio of two
back-adjusted series is not the ratio of the two contracts' prices (and a level can be negative, as in April 2020), and the
twelve-month percentage change divides a dollar change by a fictitious base. Carry must be computed from the raw contracts on
each date; the trend from ratio-adjusted returns (or from the back-adjusted change divided by the raw price held).
\end{solution}

\begin{solution}{pb:rs:market-data-for-research:1}
\textbf{1.} 1\,555\,872 messages and 55\,542 trades: 28.0 messages per trade. \textbf{2.} Additions 49.3\%, cancellations 47.2\%,
executions 3.6\%: almost every order added is later cancelled, since market orders remove only a small share of the resting
volume. \textbf{3.} $1\,555\,872 \times 36 = 56.0$ MB. \textbf{4.} Every coarser view can be rebuilt from them, and \omterm{def:rs:market-data-for-research:bar}{bars} fixed
today would freeze a clock, a threshold and a price choice into every later study. \textbf{5.} Time 390, tick 391, volume 388,
dollar 388. \textbf{6.} 5.47, 3.58, 4.04 and 3.92. \textbf{7.} 0.94; $3(1 + 0.94^2) = 5.67$, against 5.47 measured.
\textbf{8.} Volatility follows the activity level, not volume exactly (and the news burst moves the efficient price faster for a
given volume), so equal-volume \omterm{def:rs:market-data-for-research:bar}{bars} still mix variances. \textbf{9.} After the open, around the news burst and before the close,
where the activity multiplier and the U-shape put the trading. \textbf{10.} Trades 0.83, 0.58, 0.96, 0.84; mids 0.23, 0.39, 0.93,
0.85 (squared per cent). \textbf{11.} The bid--ask bounce (\cref{prop:rs:market-data-for-research:bounce}). \textbf{12.} The mid
adjusts to the efficient price in steps, so its short-horizon returns are positively autocorrelated. \textbf{13.} From about 45
seconds on (0.95 against 0.91, then within 3.4\%). \textbf{14.} A noise-robust estimator on the one-second data: two-scales
realised variance, the realised kernel or pre-averaging. \textbf{15.} \omterm{def:rs:market-data-for-research:bar}{Time bars} at five minutes or more (or a noise-robust
estimator) for the day's variance; an event clock, or the messages themselves, for short-term moves. \textbf{16.} Nothing of the
activity: it keeps a regular calendar, aligned across instruments, but its returns mix quiet and busy periods. \textbf{17.} A
common calendar: its \omterm{def:rs:market-data-for-research:bar}{bars} end at different times on different instruments and on different days. \textbf{18.} \emph{Named
result:} the one-minute \omterm{def:rs:market-data-for-research:bar}{time bars} have a kurtosis of 5.47 and the \omterm{def:rs:market-data-for-research:bar}{dollar bars} 3.92, and at one second the trade-price variance is
3.6 times the midquote's. \textbf{19.} The random activity level that scales both volume and the efficient price's speed; with
\texttt{act\_vol=0} only the U-shape and the news burst remain. \textbf{20.} None: the efficient price is unobserved, and trades,
mids and \omterm{def:rs:market-data-for-research:bar}{bars} are noisy views of it, each with its own error.
\end{solution}

\begin{solution}{iq:rs:market-data-for-research:1}
Because price variance accumulates with trading, not with time; \omterm{def:rs:market-data-for-research:bar}{bars} of equal volume (or value) have returns much closer to
Gaussian and less heteroskedastic, which helps statistics that assume both. The cost: \omterm{def:rs:market-data-for-research:bar}{bars} no longer align in time across
instruments, and the clock itself reacts to news.
\iqlookfor{subordination to volume; the loss of a common calendar.}
\end{solution}

\begin{solution}{iq:rs:market-data-for-research:2}
Microstructure noise: the bid--ask bounce adds about $s^2/2$ per return, which dominates at high frequency. The signature plot
shows it; fixes are sampling less often, using midquotes, or a noise-robust estimator.
\iqlookfor{the bounce, the signature plot, and a fix.}
\end{solution}

\begin{solution}{iq:rs:market-data-for-research:3}
Roll on a rule a real position would follow (a fixed number of days before expiry, or on open interest); compute returns from
the held contract (ratio-adjusted), and signals on levels from the contract held at the time. Back-adjusted levels rewrite the
past at every roll and can go negative; percentage changes of them are wrong.
\iqlookfor{returns versus levels, and that the adjustment is a modelling choice.}
\end{solution}

\begin{solution}{iq:rs:market-data-for-research:4}
Raw messages immutable, with exchange and receive timestamps and sequence numbers, partitioned by date and instrument,
compressed; about 76 TB a year raw at 40 bytes and 1.5 million messages per instrument-day. Derived tables (books, \omterm{def:rs:market-data-for-research:bar}{bars}) rebuilt by
versioned code, stored columnar, partitioned the way they are read, each with its lineage recorded.
\iqlookfor{immutability, lineage, the partitioning, and an order-of-magnitude size.}
\end{solution}

\begin{solution}{iq:rs:market-data-for-research:5}
Non-synchronous closes: the last trade of an illiquid stock may be minutes old, which creates spurious lead--lag and shrinks
correlations (the Epps effect); the bounce in close prices; heteroskedasticity from uneven activity; auction prints in the first
and last \omterm{def:rs:market-data-for-research:bar}{bars}.
\iqlookfor{non-synchronicity and noise, not only look-ahead.}
\end{solution}

\begin{solution}{iq:rs:market-data-for-research:6}
Given $V$, $r \sim \mathcal N(0, \sigma^2V)$, so $\E[r^2] = \sigma^2\E[V]$ and $\E[r^4] = 3\sigma^4\E[V^2]$; the ratio is $3\E[V^2]/\E[V]^2
= 3(1 + \mathrm{CV}^2)$.
\iqlookfor{conditioning on the volume and the fourth moment of a normal.}
\end{solution}
